\section{The interpretations, including both composites}
\label{sec:interpretations}

The equivalence in Theorem~\ref{rev:equivalence} is stronger than agreement
on standard numerical computations: it holds for every sequent, with
arbitrary retained assumptions. This is what permits reuse of the
existing bi-interpretation rather than only an equiconsistency argument.
For clarity we spell out the translations and the role of the inherited
arithmetic results.

\subsection{Which iPRA is intended}

Here iPRA is full first-order intuitionistic logic with equality, a
function symbol for every fixed primitive-recursive definition, its
defining equations, the successor axioms, and induction for
quantifier-free formulas. Quantified formulas are allowed everywhere
in the logic; induction is not extended to arbitrary quantified formulas.
The pinned \code{PR.lean} implements this convention~\cite{leanbridge}.

The binding notation $I(n,a;u.h)$ denotes a fixed finite description.
It can be elaborated into a staged first-order signature: for each
earlier finite body $h(u,\vec p)$ introduce $i_h(n,a,\vec p)$ with
the two iterator equations. Parameters in the description become
ordinary first-order arguments. No variable ranges over functions or
programs, and no uniformly total evaluator for all primitive-recursive
descriptions is asserted.

\subsection{Native formulas and ordinary intuitionistic logic}

Let $T_I$ be ordinary intuitionistic first-order logic over the scalar
terms, with the two computation equations, induction for literal
equations, coverage
\[
  \forall x\,(x=0\vee\exists y\,x=Sy),
\]
and the explicit separation bridge $(0=S0)\to\bot$. Translate native
equations identically, and translate $Q$ and $D$ by the meanings given
in Section~\ref{rev:section}. Both interpretations use the full domain,
defined by $x=x$, and interpret equality as equality.

The reverse logical translation uses only native formulas. Put
$C(v,a,b)=I(v,a;u.b)$, with the state variable unused. Its values at zero
and successors are $a$ and $b$. The gate $G_vA$ is true at zero and
equivalent to $A$ at a successor; the selector $M_v(A,B)$ selects $A$
at zero and $B$ at successors. Both are constructed by external
recursion on fixed finite formulas, using equation selectors and
padding to align binder depth. Their full definitions and proof rules
are reproduced in Appendix~\ref{sec:choice-introduction}.

With $p$ the parity iterator, one convenient Eq/$D$ macro inventory is
\[
\begin{aligned}
 \bot_*&=(0=S0),\\
 A\to_*B&=D_{u,z}M_z(A,B),\\
 A\land_*B&=D_{u,z}G_zM_{p(z)}(A,B),\\
 \exists_*x A(x)&=D_{x,y}G_yA(x),\\
 A\lor_*B&=\exists_*x M_x(A,B),\\
 \forall_*x A(x)&=A(0)\land_*D_{u,z}G_zA(z).
\end{aligned}
\]
Here the unused coordinate $u$ and all new bound variables are chosen
fresh. These are definitions, not new axiom schemas. The logical
introduction and elimination rules, falsity elimination, numerical
coverage, and both native/ordinary formula composites follow by the
explicit constructions in the appendix. An auxiliary $Q_yB$ is
interderivable with $D_{u,y}B$ when $u$ is unused.

The ordinary separation bridge must not be omitted: positive numerical
equations alone admit a one-element model. Native falsity is a designated
equation, and the bridge identifies its intended logical role in $T_I$.
This is already part of the inherited presentation, not an additional
assumption introduced for reversal.

\subsection{Scalar arithmetic and primitive recursion}

Interpret an iterator in iPRA by primitive recursion whose update ignores
the recursion counter. This translates the scalar equations and literal
equation induction directly. Coverage and separation are iPRA theorems.

For the converse translation, the inherited internal derivations proceed
in the following dependency order.
\begin{enumerate}
\item A scalar predecessor $P$ is constructed and proved to satisfy
  $P0=0$ and $P(Sn)=n$ using equation induction and numerical cases.
  Pairing and successor injectivity are not assumed in that bootstrap.
\item Truncated subtraction and a Boolean reflection of actual equality
  are derived. Each fixed quantifier-free formula is represented by a
  characteristic equation, giving its induction instance.
\item With quantifier-free induction available, pairing $\pi$ and
  projections $L,H$ are constructed internally. They satisfy
  $L(\pi(a,b))=a$ and $H(\pi(a,b))=b$ throughout the domain.
\item Each fixed primitive recursion is compiled by storing its counter
  and current value in one paired scalar state.
\end{enumerate}

For example, for
\[
 f(0,\vec p)=g(\vec p),\qquad
 f(Sn,\vec p)=h(n,f(n,\vec p),\vec p),
\]
the translation uses
\[
\begin{split}
 V(n,\vec p)&=I\bigl(n,\pi(0,g^*(\vec p));\,
 u.\pi(SL(u),h^*(L(u),H(u),\vec p))\bigr),\\
 f^*(n,\vec p)&=H(V(n,\vec p)).
\end{split}
\]
Equation induction proves $L(V(n,\vec p))=n$, and the projection
identities establish the defining equations of $f$. Translation is
external recursion on a fixed finite definition, not an internal
universal-recursion principle. The complete inherited constructions
appear in the scalar predecessor, reflection, and primitive-recursion
appendices.

\subsection{Why both round trips are identities}

For a fixed iterator, equation induction proves uniqueness from its
zero and successor equations. The analogous uniqueness proof applies
to each fixed primitive recursion. Inducting externally over a finite
function description therefore proves that its double translation is
equal to the original function inside the corresponding theory. The
interpretation domain is never changed, so the graph of the resulting
domain isomorphism is literally $x=y$.

For formulas, use the proved meanings of the logical macros and the
native $Q,D$ binders, then induct on formula structure. For example,
the native paired composite is
\[
 D_{x,y}B\ \dashv\vdash\
 \exists_*x\bigl(B(x,0)\to_*\forall_*y B(x,y)\bigr).
\]
The forward direction opens a fresh witness and introduces the three
logical macros; the backward direction eliminates the existential
macro and uses DI-case with that witness. The appendix gives both
directions with their local assumptions.

Combining these arithmetic and logical composites with
Theorem~\ref{rev:equivalence} establishes the requested bi-interpretation
for $\Rcal$. In particular,
\[
 \varnothing\vdash_{\Rcal}0=S0
 \quad\Longleftrightarrow\quad
 \ipra\vdash\bot.
\]
This is a relative consistency equivalence; no assertion of the
consistency of either theory is needed as a premise.

\subsection{Direct closure versus the implemented progress packet}
\label{sec:progress-packet}

The mathematical presentation in Section~\ref{rev:section} follows the
current Lean kernel and introduces $Q_yE(y)$ directly by equation
progress. The tiny inherited certificate format instead introduces
$D_{x,y}E(x,y)$ from $E(t,z)\vdash E(t,Sz)$, with a witness $t$.
These are equivalent interfaces, not literally the same instruction.

To obtain witnessed progress from direct closure, set $F(y)=E(t,y)$.
QI-progress gives $Q_yF(y)$. Under $F(0)$, ordinary QE gives $F(Sz)$;
DI-case with witness $t$ finishes. Ordinary QE is already derived in
$\Rcal$ by reversing QX's result.

Conversely, apply witnessed progress with an unused witness coordinate
$u$ to obtain $D_{u,y}E(y)$. Eliminate it. The supplied assumption is
exactly $Q_yE(y)$, independent of the fresh witness, so hypothesis
closes the branch. Both constructions are valid in arbitrary retained
contexts and keep the progress matrix a literal equation. The supplied
producer uses the witnessed packet convention; the Lean draft uses the
direct-closure convention.
